\section{Conclusion}

Decades after its introduction, web tracking still remains an \textit{archetypal cat-and-mouse game}.
%
Each incremental defense---whether a new browser policy or a regulatory ruling---quickly provokes an equally sophisticated evasion technique to track users.
%
This adversarial dynamic shows that purely reactive approaches cannot deliver privacy guarantees for online users.
%
\textit{Regulations alone are insufficient}; data protection statutes such as GDPR and CCPA have tightened accountability, yet enforcement lags the speed of technical changes in evolving tracking mechanisms.
%
Moreover, trackers often find tolerated gray zones to bypass regulations.
%
As a result, enforcement frequently stalls on jurisdictional or interpretative disputes.
%
Regulators and the measurement community need to collaborate to incorporate agile, evidence-driven auditing methods to avoid any oversight and ensure that regulations evolve competitively with the technical reality.
%
On the other hand, while \textit{browsers can be powerful gatekeepers, they provide an unreliable line of defense}.
%
Default protections vary widely across browser vendors, mitigations ship years after the issues are identified, and commercial incentives often result in more permissive designs.
%
Future research must therefore look beyond ``\textit{fix it in the browser}'' remedies and explore complementary approaches that truly safeguard users' privacy.

Thus, while browser-based protections and policy-driven changes are effective to some extent, the current tracking landscape demands a \textbf{default privacy-first solution} where users can control their privacy as opposed to browsers or regulators.
%
By systematizing evolution of web tracking, its defenses, and regulatory compliance, this SoK suggests key future directions and calls for a privacy-first web in which tracking is systematically measured and proactively audited, while users remain in control of their data by default.
